\chapter[ESTADOS DO MODELO DO TEP]{ESTADOS DO MODELO DO TEP}
\label{apendice:estados}

Este apêndice lista as 50 variáveis de estado do modelo de \citeonline{DOWNS1993}, na ordem do vetor \texttt{YY} do código FORTRAN original (sub-rotina \texttt{TEFUNC}), e a chave com que cada uma aparece no \texttt{tep-plant}. A Tabela~\ref{tab:estados_tep} resume os estados por grupo.

No reator e no separador, os componentes A, B e C, que não condensam, são guardados só como vapor, e os componentes D a H só como líquido. No \textit{stripper} só há líquido, e no vaso do compressor só há vapor. Cada equipamento tem ainda um estado de energia interna. Os estados 37 e 38 são as temperaturas de saída da água de resfriamento do reator e do condensador, e os 12 últimos são as posições das válvulas, que em regime coincidem com as 12 XMV (Seção~\ref{subsec:otimizacao}).

\begingroup
\small
\begin{longtable}{|c|l|>{\raggedright\arraybackslash}p{4.3cm}|>{\raggedright\arraybackslash}p{6.0cm}|}
	\caption{Variáveis de estado do modelo do TEP e chaves correspondentes no \texttt{tep-plant}.}
	\label{tab:estados_tep} \\
	\hline
	\textbf{Estado} & \textbf{FORTRAN} & \textbf{Descrição} & \textbf{Chave no \texttt{tep-plant}} \\
	\hline
	\endfirsthead
	\hline
	\textbf{Estado} & \textbf{FORTRAN} & \textbf{Descrição} & \textbf{Chave no \texttt{tep-plant}} \\
	\hline
	\endhead
	\hline
	\endfoot
	1--3   & \texttt{UCVR(1--3)} & Quantidade de A, B e C no vapor do reator (kmol) & \texttt{reactor.\allowbreak{}state.\allowbreak{}vapor\_a} a \texttt{vapor\_c} \\
	4--8   & \texttt{UCLR(4--8)} & Quantidade de D a H no líquido do reator (kmol) & \texttt{reactor.\allowbreak{}state.\allowbreak{}liquid\_d} a \texttt{liquid\_h} \\
	9      & \texttt{ETR}        & Energia interna do reator & \texttt{reactor.\allowbreak{}state.\allowbreak{}enthalpy} \\
	\hline
	10--12 & \texttt{UCVS(1--3)} & Quantidade de A, B e C no vapor do separador (kmol) & \texttt{separator.\allowbreak{}state.\allowbreak{}vapor\_a} a \texttt{vapor\_c} \\
	13--17 & \texttt{UCLS(4--8)} & Quantidade de D a H no líquido do separador (kmol) & \texttt{separator.\allowbreak{}state.\allowbreak{}liquid\_d} a \texttt{liquid\_h} \\
	18     & \texttt{ETS}        & Energia interna do separador & \texttt{separator.\allowbreak{}state.\allowbreak{}enthalpy} \\
	\hline
	19--26 & \texttt{UCLC(1--8)} & Quantidade de A a H no líquido do \textit{stripper} (kmol) & \texttt{stripper.\allowbreak{}state.\allowbreak{}0} a \texttt{7} \\
	27     & \texttt{ETC}        & Energia interna do \textit{stripper} & \texttt{stripper.\allowbreak{}state.\allowbreak{}8} \\
	\hline
	28--35 & \texttt{UCVV(1--8)} & Quantidade de A a H no vapor do vaso do compressor (kmol) & \texttt{compressor.\allowbreak{}state.\allowbreak{}0} a \texttt{7} \\
	36     & \texttt{ETV}        & Energia interna do vaso do compressor & \texttt{compressor.\allowbreak{}state.\allowbreak{}8} \\
	\hline
	37     & \texttt{TWR}        & Temperatura de saída da água de resfriamento do reator (°C) & constante: \texttt{heat.\allowbreak{}reactor\_cooling\_water\_return} \\
	38     & \texttt{TWS}        & Temperatura de saída da água de resfriamento do condensador (°C) & constante: \texttt{heat.\allowbreak{}separator\_cooling\_water\_return} \\
	\hline
	39     & \texttt{VPOS(1)}    & Válvula de alimentação de D, XMV$_1$ (\%) & \texttt{valve.\allowbreak{}feed\_d.\allowbreak{}position} \\
	40     & \texttt{VPOS(2)}    & Válvula de alimentação de E, XMV$_2$ (\%) & \texttt{valve.\allowbreak{}feed\_e.\allowbreak{}position} \\
	41     & \texttt{VPOS(3)}    & Válvula de alimentação de A, XMV$_3$ (\%) & \texttt{valve.\allowbreak{}feed\_a.\allowbreak{}position} \\
	42     & \texttt{VPOS(4)}    & Válvula de alimentação de A e C, XMV$_4$ (\%) & \texttt{valve.\allowbreak{}feed\_ac.\allowbreak{}position} \\
	43     & \texttt{VPOS(5)}    & Válvula de reciclo do compressor, XMV$_5$ (\%) & \texttt{valve.\allowbreak{}compressor\_recycle.\allowbreak{}position} \\
	44     & \texttt{VPOS(6)}    & Válvula de purga, XMV$_6$ (\%) & \texttt{valve.\allowbreak{}purge.\allowbreak{}position} \\
	45     & \texttt{VPOS(7)}    & Válvula de líquido do separador, XMV$_7$ (\%) & \texttt{valve.\allowbreak{}separator\_underflow.\allowbreak{}position} \\
	46     & \texttt{VPOS(8)}    & Válvula de produto do \textit{stripper}, XMV$_8$ (\%) & \texttt{valve.\allowbreak{}stripper\_product.\allowbreak{}position} \\
	47     & \texttt{VPOS(9)}    & Válvula de vapor do \textit{stripper}, XMV$_9$ (\%) & \texttt{valve.\allowbreak{}stripper\_steam.\allowbreak{}position} \\
	48     & \texttt{VPOS(10)}   & Água de resfriamento do reator, XMV$_{10}$ (\%) & \texttt{valve.\allowbreak{}reactor\_cooling\_water.\allowbreak{}position} \\
	49     & \texttt{VPOS(11)}   & Água de resfriamento do condensador, XMV$_{11}$ (\%) & \texttt{valve.\allowbreak{}condenser\_cooling\_water.\allowbreak{}position} \\
	50     & \texttt{VPOS(12)}   & Velocidade do agitador, XMV$_{12}$ (\%) & \texttt{agitator.\allowbreak{}speed} \\
\end{longtable}
\endgroup
\vspace{-1.5em}
\fonte{Elaborado pelo autor, com base no código de \citeonline{DOWNS1993} e no \texttt{tep-plant}.}

No \texttt{tep-plant}, os estados 37 e 38 não são integrados: as duas temperaturas ficam fixas nos valores do caso base, 94,6~°C no reator e 77,3~°C no condensador. As chaves da tabela são as que o simulador publica para leitura; os nomes de configuração do caso base, usados para carregar os valores iniciais, estão no arquivo \texttt{te\_mode1\_initial\_state.\allowbreak{}toml}.
